\section{Introduction}
\label{sec:introduction}

For a complex Hilbert space $H$, let $K(H)$ denote the $C^*$-algebra of
compact operators on $H$. Every nonzero irreducible representation of
$K(H)$ is unitarily equivalent to its defining representation.
Naimark's problem asks whether this property characterizes the algebras
$K(H)$. We give a negative answer in ZFC.

\begin{theorem}\label{na:main}
There exists a unital infinite-dimensional simple complex $C^*$-algebra
$A$ with a faithful tracial state such that all its nonzero irreducible
representations are unitarily equivalent. In particular, $A$ is not
isomorphic to $K(H)$ for any complex Hilbert space $H$.
\end{theorem}

Here a representation is irreducible if it is nonzero and has no
nontrivial closed invariant subspace. A tracial state $\tau$ is faithful if
$\tau(a^*a)>0$ whenever $a\ne0$. The last assertion of the theorem follows
from the first: $K(H)$ is unital only when $H$ is finite-dimensional.
All algebras and Hilbert spaces in the proof are complex, and no
separability assumption is placed on $A$.

\subsection*{History and scope}
Naimark's question dates to 1951; see the historical account of the
question and Rosenberg's affirmative work in \cite{AkemannWeaver2004}.
When the unique nonzero irreducible representation acts on a separable
Hilbert space, the algebra is isomorphic to the compact operators;
see \cite[Proposition~7.2]{CalderonFarah2022}. This condition concerns
an irreducible representation, not merely a faithful representation
on a separable Hilbert space.
Akemann and Weaver constructed a counterexample
assuming Jensen's diamond principle
\cite[Theorem~5]{AkemannWeaver2004}. They also established independence
for the existence of a counterexample generated by $\aleph_1$ elements
\cite[Corollary~7]{AkemannWeaver2004}. The unrestricted existence question
does not impose this bound on the number of generators.

Subsequent work refined the set-theoretic hypotheses and the possible
structure of counterexamples. Calder\'on and Farah obtain a
counterexample from $\diamondsuit_{\mathrm{Cohen}}$ together with CH
\cite[Theorem~5.4]{CalderonFarah2022}. Under diamond, Vaccaro constructs
counterexamples whose tracial state spaces are affinely homeomorphic
to any prescribed metrizable Choquet simplex
\cite{Vaccaro2018,Vaccaro2020}. Thus the presence of a trace is already
compatible with the earlier constructions.

Tanaka gives another ZFC construction with a faithful tracial state
\cite[Theorems~1.1, 4.1 and~5.3]{Tanaka2026}, so the existence conclusion
of Theorem~\ref{na:main} is already present there. His construction
starts from the CAR algebra and uses strong shell sums and rank-one
defects. It also gives a faithful representation on a separable Hilbert
space; that representation is not irreducible. Here we give an
alternative proof organized around a selective bridge for faithfully
tracial bases without finite-dimensional irreducible representations,
reduced-HNN reordering, and dependency-closed
unique extension and countable determination. The precise hypotheses
and conclusions of these intermediate results are stated below.
Our construction has length $\mathfrak c^+$, where
$\mathfrak c=2^{\aleph_0}$; we make no assertion that this length is the
density of the resulting algebra.

\subsection*{The proof in outline}
We work with pure states, whose Gelfand--Naimark--Segal (GNS)
representations are irreducible.
Two pure states are called equivalent when their GNS representations
are unitarily equivalent. The construction repeatedly joins a chosen
pair of equivalence classes. For this procedure to be controllable,
each old pure state must have a unique state extension, and no other
pair of old classes may become equivalent.

The main analytic step starts with a faithfully tracial algebra having
no finite-dimensional irreducible representations. It adjoins a unitary
carrying one decreasing projection sequence to another. Each sequence specifies its state as
the only state taking value one on every projection; we call it a
\emph{peaking sequence}. When the two states are inequivalent and the
sequences have matching trace
values strictly decreasing to zero, the resulting reduced HNN extension
has exactly the required pure-state properties
(Theorem~\ref{na:bridge}). The extension is built from a reduced
amalgamated free product using Ueda's matrix-corner construction
\cite{UedaPreprint2008,Ueda2008}. Pure-state excision
\cite{AkemannAndersonPedersen1986} and the cross-excision estimate of
Akemann and Weaver \cite[Lemma~1]{AkemannWeaver2004} turn uniqueness
of extension and the exclusion of unintended joins into norm-decay
estimates for reduced words. The estimates apply to every old pure state.

The second step makes every pure state accessible to the recursion.
We start with a separable simple algebra, adjoin one unitary at each
active successor stage, and take norm closures of unions at limits.
Each adjunction depends on only countably many earlier generators.
If a set of generators contains all these dependencies, we prove that
every pure state on the algebra it generates has a unique extension
to the whole algebra (Theorem~\ref{na:closed-extension}). To prove this,
we postpone an omitted adjunction past the retained ones. The
no-other-join assertion of the bridge keeps its two states inequivalent
until the postponed adjunction is performed. For every pure state of
the full algebra, excision provides a separable subalgebra generated
by a countable dependency-closed set of steps on which the restriction
is pure. The closed-subiteration theorem then makes the original state
its unique extension to the full algebra
(Lemma~\ref{na:countable-determination}).

This last conclusion replaces the state-guessing role of a
set-theoretic prediction principle. At every stage below
$\mathfrak c^+$ there are at most $\mathfrak c$ pure states, so a fixed
schedule can eventually treat them all. On each separable simple
determining subalgebra, the homogeneity theorem of Kishimoto, Ozawa and
Sakai \cite{KishimotoOzawaSakai2003} supplies the next peaking sequence.
A final countable-determination argument catches every pure state of
the limit algebra.

The ingredients have several useful precedents. Akemann and Weaver
combine pure-state homogeneity with extensions joining prescribed
states \cite[Corollary~3]{AkemannWeaver2004}. Farah and Hirshberg realize prescribed countable equivalence patterns
with unique pure-state extensions
\cite[Lemma~2.3]{FarahHirshberg2017}; Vaccaro's corrected construction
combines this method with control of the trace space
\cite{Vaccaro2018,Vaccaro2020}.
Akemann, Wassermann and
Weaver prove pure-state restriction and unique-extension results for
reduced free group algebras
\cite[Proposition~2.3 and Theorem~2.6]{AkemannWassermannWeaver2010}.
Calder\'on and Farah propose using separable subalgebras that uniquely
determine the relevant states, followed by an amalgamation preserving
unique extension \cite[Section~8, especially Question~8.1]{CalderonFarah2022}.
The two features established here are the selective bridge for all
old pure states and its compatibility with dependency-closed subiterations.
Together they give the countable determination needed for the
unrestricted recursion.

Section~\ref{sec:states} proves the state-theoretic facts used below.
Section~\ref{sec:hnn} constructs the reduced adjunction, and
Section~\ref{sec:bridge} proves its pure-state properties.
Section~\ref{sec:iteration} establishes countable determination;
Section~\ref{sec:construction} carries out the final recursion.
